
\SGAProof{Proof}
Parts \SGARefLink{sga3:xvii:appendix:A}{(a)}, \SGARefLink{sga3:xvii:section:b}{(b)}, \SGARefLink{sga3:xvii:appendix:C}{(c)}, and (d) follow from \SGARefLink{sga3:xvii:proposition:5:3:1}{5.3.1}.

\smallskip
\noindent
(e) Let \(Z\) be the centralizer of \(H'\) in \(U\), and let \(T\)
be the transporter of \(H'\) into \(H''\).  By \SGARefLink{sga3:xvii:appendix:A}{(a)}, \(T\) is
nonempty; hence \(T\) is a principal homogeneous space under \(Z\).
The hypothesis on \(Z\) therefore implies that \(T\) is trivial.

\smallskip
\noindent
(f) Proceeding as in \SGARefLink{sga3:xvii:proposition:5:3:1}{5.3.1}, we see that it is enough to consider the
case where \(H\) is étale.  Suppose first that \(H\) is
diagonalizable, defined by an abstract group \(M\) of order \(m\)
prime to \(p\).  The two lifts \(H'\) and \(H''\) determine a
one-cocycle of \(M\) with values in \(U(k)\), that is, a map
\[
  h:M\longrightarrow U(k)
\]
such that
\[
  h(mn)=h(m)\,{}^{m}\!h(n)
\]
for every pair \(m,n\) of elements of \(M\).  The groups \(H'\) and
\(H''\) are conjugate by an element of \(U(k)\) if and only if there
exists \(a\in U(k)\) such that
\[
  h(m)=a^{-1}({}^{m}\!a).
\]
Now the abstract group \(U(k)\) has a composition series whose
successive quotients are commutative and annihilated by a power of
\(p\) (in view of \SGARefLink{sga3:xvii:proposition:5:6:1:item:c}{5.6.1(c)}, we may suppose that \(p>0\)).  It follows
at once in this case that \(h\) is a coboundary.

Consider now the general case.  Again let \(Z\) denote the
centralizer of \(H'\) in \(U\), and \(T\) the transporter of \(H'\)
into \(H''\); thus \(T\) is a torsor under \(Z\).  By hypothesis there
is a finite Galois extension \(K/k\), with Galois group \(G\) of
order prime to \(p\), which splits \(H\).  By the preceding case,
\(H'_K\) and \(H''_K\) are conjugate by an element of \(U(K)\), and
therefore \(T_K\) is trivial.  Consequently \(T\) is defined by an
element of
\[
  H^1(G,Z).
\]
For the same reasons as above, the hypothesis on \(G\) implies that
\[
  H^1(G,Z)=e;
\]
hence \(T\) is trivial.

\SGASubsection{5.7}{Proof of \SGARefLink{sga3:xvii:theorem:5:1:1:item:ii}{5.1.1(ii)}\SGARefLink{sga3:xvii:section:b}{(b)}}
\SGARefTarget{sga3:xvii:subsection:5:7}

\SGAStatement{Lemma 5.7.1}{Let \(S\) be a prescheme, let \(G\) be an
\(S\)-group prescheme, separated and smooth over \(S\), and let \(H\)
be an \(S\)-group of multiplicative type acting on \(G\).  Then the
\(S\)-functor
\[
  Z=G^H
\]
of invariants of \(G\) under \(H\) is representable by a closed
smooth \(S\)-group subprescheme of \(G\).}
\SGARefTarget{sga3:xvii:lemma:5:7:1}

The representability of \(G^H\) by a closed group subprescheme of
\(G\) follows from \SGARefLink{sga3:viii:examples:6:5:item:d:01}{Exposé VIII, 6.5(d)}.  Consider the semidirect
product
\[
  K=G\rtimes_S H.
\]
The centralizer of \(H\) in \(K\) is then \(Z\times_S H\).  To prove
that \(Z\) is smooth, we must verify the following.  Suppose that
\(S\) is affine, that \(S_0\) is a closed subscheme of \(S\) defined
by a square-zero ideal, and that \(u_0\in Z(S_0)\).  We must lift
\(u_0\) to an element of \(Z(S)\).

Since \(G\) is smooth over \(S\), there is an element \(u\in G(S)\)
lifting \(u_0\).  Let \(i\) be the canonical immersion of \(H\) into
\(K\), and consider the two \(S\)-group morphisms
\[
  H\rightrightarrows K,\qquad i
  \quad\hbox{and}\quad
  j=\operatorname{int}(u)i.
\]
Since \(u_0\in Z(S_0)\), we have \(i_{S_0}=j_{S_0}\).  By
\SGARefLink{sga3:ix:theorem:3:2}{Exposé IX, 3.2}, there exists \(v\in K(S)\), lifting the identity
section in \(K(S_0)\), such that
\[
  i=\operatorname{int}(v)j
   =\operatorname{int}(vu)i.
\]
Thus
\[
  vu=(u',v')\in (Z\times_S H)(S).
\]
It follows that \(u'\in Z(S)\), and \(u'\) lifts \(u_0\).

\SGAStatement{Lemma 5.7.2}{Let \(k\) be a field and \(H\) a
\(k\)-algebraic group, and let \(P(H)\) denote the following
property:
\begin{quote}
For every smooth unipotent \(k\)-group \(U\), and every extension
\(E\) of \(H\) by \(U\), any two lifts of \(H\) in \(E\) are
conjugate by an element of \(U(k)\).
\end{quote}
Suppose that \(H\) is an algebraic group which is an extension of a
group \(H''\) of multiplicative type by a group \(H'\) of
multiplicative type.  If \(P(H')\) and \(P(H'')\) hold, then
\(P(H)\) holds.}
\SGARefTarget{sga3:xvii:lemma:5:7:2}

Indeed, let \(U\) be a smooth unipotent \(k\)-group, let \(E\) be an
extension of \(H\) by \(U\), and let \(H_1,H_2\) be two lifts of
\(H\) in \(E\).  Let \(H'_1,H'_2\) be the corresponding lifts of
\(H'\).  Since \(P(H')\) holds, there exists \(u\in U(k)\) such that
\[
  H'_2=\operatorname{int}(u)H'_1.
\]
After replacing \(H_1\) by \(\operatorname{int}(u)H_1\), we may
suppose that \(H'_1=H'_2\); denote this common group simply by
\(H'\).  Put
\[
  E'=\operatorname{Centr}_E(H')
     =U^{H'}H_1
     =U^{H'}H_2.
\]
By \SGARefLink{sga3:xvii:lemma:5:7:1}{5.7.1}, \(U^{H'}=U'\) is smooth.  We then have an extension
\[
  1\longrightarrow U'\longrightarrow E'
    \longrightarrow H\longrightarrow1.
\]
By construction \(H'\) is central in \(E'\), and hence invariant.
Passing to the quotient gives an exact sequence
\[
  1\longrightarrow U'\longrightarrow E''
    \longrightarrow H''\longrightarrow1.
\]
Since \(U'\) is smooth and \(P(H'')\) holds, the images of \(H_1\)
and \(H_2\) in \(E''\) are conjugate by an element
\(u\in U'(k)\).  It follows that
\[
  H_1=\operatorname{int}(u)H_2.
\]

To prove \SGARefLink{sga3:xvii:theorem:5:1:1:item:ii}{5.1.1(ii)}\SGARefLink{sga3:xvii:section:b}{(b)}, observe that, since \(k\) is algebraically
closed, \(H\) has a composition series whose successive quotients
are smooth or, when \(p>0\), isomorphic to \(\boldsymbol\mu_p\).
Repeated application of \SGARefLink{sga3:xvii:lemma:5:7:2}{5.7.2} reduces us to the case where \(H\) is
smooth or \(H=\boldsymbol\mu_p\).  In the first case it is enough to
apply \SGARefLink{sga3:xvii:theorem:5:1:1:item:ii}{5.1.1(ii)}(a).  It remains to consider \(H=\boldsymbol\mu_p\).
Since \(U\) is smooth, \(U\) has a characteristic composition series
whose successive quotients are étale or isomorphic to
\(\mathbb G_a^r\), by \SGARefLink{sga3:xvii:lemma:3-9}{3.9}.  If \(U\) is étale, apply \SGARefLink{sga3:xvii:lemma:5:5:1}{5.5.1}.  We are
finally reduced to
\[
  H=\boldsymbol\mu_p,\qquad U=\mathbb G_a^r.
\]

We must prove that
\[
  H^1(\boldsymbol\mu_p,U)=0.
\]
The method used in \SGARefLink{sga3:xvii:proposition:5:4:1}{5.4.1} no longer applies, because
\(\boldsymbol\mu_p\) does not in general act linearly on
\(\mathbb G_a^r\).  Fix the following notation: \(H'\) denotes a
lift of \(H\) in \(E\), and
\[
  \mathfrak e=\operatorname{Lie}E,\qquad
  \mathfrak u=\operatorname{Lie}U,\qquad
  \mathfrak h=\operatorname{Lie}H,\qquad
  \mathfrak h'=\operatorname{Lie}H'.
\]
Let \(X\) be a nonzero element of \(\mathfrak h\) such that
\[
  X^{(p)}=X
\]
(Appendix II, \SGARefLink{sga3:xvii:lemma:3-1}{3.1}), and let \(X'\) be its lift in \(\mathfrak h'\).

Since \(\boldsymbol\mu_p\) is a radicial group of height \(1\), the
set of lifts of \(H\) in \(E\) is in bijective correspondence with
the set \(A\) of elements \(Y\in\mathfrak e\) such that
\[
  Y^{(p)}=Y
\]
and whose image is \(X\) (Appendix II, \SGARefLink{sga3:ii:subsection:2-2}{2.2}).  Likewise, if
\(Y\in A\) corresponds to a lift \(H''\) of \(H\) in \(E\), and if
\(u\in U(k)\), then
\[
  \operatorname{int}(u)H'=H''
  \quad\Longleftrightarrow\quad
  \operatorname{Ad}(u)X'=Y.
\]
Let \(B\) be the subset of \(\mathfrak e\) consisting of the
elements \(\operatorname{Ad}(u)X'\), with \(u\in U(k)\).  Plainly
\(B\subset A\), and it remains only to prove \(A=B\) when \(k\) is
algebraically closed.

\smallskip
\noindent
(a) \emph{Study of \(A\).}
Since \(\mathfrak u\) is commutative and normalized by
\(\mathfrak h'\), Jacobson's formula (Exposé VII\textsubscript{A},
\SGARefLink{sga3:VIIB:5.2:theorem}{5.2}) gives, for \(u\in\mathfrak u\),
\[
  (X'+u)^{(p)}
   =X'{}^{(p)}+u^{(p)}
      +(\operatorname{ad}X')^{p-1}(u)
   =X'+(\operatorname{ad}X')^{p-1}(u).
\]
Consequently,
\[
  X'+u\in A
  \quad\Longleftrightarrow\quad
  u=(\operatorname{ad}X')^{p-1}(u).
\]
Let
\[
  \mathfrak u
   =\bigoplus_{n\in\mathbb Z/p\mathbb Z}\mathfrak u_n
   =\mathfrak u_0\oplus
      \bigoplus_{n\in(\mathbb Z/p\mathbb Z)^\times}\mathfrak u_n
\]
be the canonical decomposition of \(\mathfrak u\) under the action
of \(H'\).  If \(u\in\mathfrak u_0\), then
\(\operatorname{ad}X'(u)=0\).  If
\(u\in\mathfrak u_n\), with
\(n\in(\mathbb Z/p\mathbb Z)^\times\), then
\[
  (\operatorname{ad}X')^{p-1}(u)=u.
\]
Thus \(Y=X'+u\) belongs to \(A\) if and only if
\[
  u\in
  \bigoplus_{n\in(\mathbb Z/p\mathbb Z)^\times}\mathfrak u_n.
\]
In particular, \(A\) is the set of \(k\)-rational points of an
irreducible subscheme of \(\mathbf W(\mathfrak e)\), of dimension
\[
  \operatorname{rk}\mathfrak u-\operatorname{rk}\mathfrak u_0
  =
  \operatorname{rk}\mathfrak u
  -\operatorname{rk}\operatorname{Centr}_{\mathfrak u}(X').
\]

\smallskip
\noindent
(b) \emph{Study of \(B\).}
We shall need the following lemma.

\SGAStatement{Lemma 5.7.3 (Rosenlicht)}{Let \(U\) be a unipotent
algebraic group over a field \(k\), acting on a quasi-affine
\(k\)-scheme \(X\).  Then the orbit of every point \(x\in X(k)\) is
closed in \(X\).  Here, by the orbit of \(x\), we mean the subset of
\(\operatorname{ens}(X)\) which is the image of \(U\) under the
morphism \(g\mapsto g\cdot x\).}
\SGARefTarget{sga3:xvii:lemma:5:7:3}

By functoriality, \(U\) acts on the affine envelope of \(X\), that
is,
\[
  \operatorname{Spec}\Gamma(X,\mathcal O_X),
\]
so we may suppose that \(X\) is affine.  We may then suppose that
\(k\) is algebraically closed, \(X\) is reduced, and \(U\) is smooth
(observe that \(U_{\mathrm{red}}\) acts on \(X_{\mathrm{red}}\) when
\(k\) is perfect).  Let \(Y\) be the schematic image of \(U\)
(EGA I, 9.5.1) under the morphism \(g\mapsto g\cdot x\).  It is a
closed reduced subscheme of \(X\), on which \(U\) acts.  It follows
readily from EGA IV, 1.8.6 that the orbit of \(x\) is an open subset
\(Z\) of \(Y\), dense in \(Y\).  We must prove that
\[
  Z=\operatorname{ens}(Y).
\]
Let \(F\) be the closed reduced subscheme of \(Y\) whose underlying
space is \(Y\setminus Z\).  Thus \(F=V(J)\), where \(J\) is a
nonzero ideal of \(\Gamma(Y,\mathcal O_Y)\).  Since \(U\) is smooth,
\(U\) acts on \(F\), hence on \(J\), and therefore
\[
  J^U\ne0
\]
by \SGARefLink{sga3:xvii:proposition:3:2}{3.2}.  If \(a\) is a nonzero element of \(J^U\), then \(a\) is
constant on the orbit \(Z\), and hence on \(Y\), because \(Z\) is
dense in \(Y\).  The ideal \(J\) consequently contains \(k\), and
\(F=\varnothing\).

Apply the preceding lemma to the group \(U\) acting on the affine
space \(\mathbf W(\mathfrak e)\) through the adjoint representation.
The orbit of \(X'\) is the underlying set of a closed subscheme of
\(\mathbf W(\mathfrak e)\).  Moreover, the stabilizer \(Z\) of
\(X'\) is its centralizer in \(U\), and there is a closed immersion
\[
  U/Z\longrightarrow\mathbf W(\mathfrak e).
\]
Thus the orbit of \(X'\) is the underlying space of a closed
subscheme of \(\mathbf W(\mathfrak e)\) of dimension
\[
  \dim U-\dim Z.
\]

\smallskip
\noindent
\SGARefLink{sga3:xvii:appendix:C}{(c)} \emph{End of the proof of \SGARefLink{sga3:xvii:theorem:5:1:1:item:ii}{5.1.1(ii)}\SGARefLink{sga3:xvii:section:b}{(b)}.}
When \(k\) is algebraically closed, the canonical map
\[
  U(k)\longrightarrow(U/Z)(k)
\]
is surjective.  By \SGARefLink{sga3:xvii:section:b}{(b)}, \(B\) is therefore the set of \(k\)-rational
points of a closed subscheme of \(\mathbf W(\mathfrak e)\) of
dimension \(\dim U-\dim Z\).  In view of \SGARefLink{sga3:xvii:appendix:A}{(a)}, to prove \(A=B\) it is
enough to show that
\[
  \operatorname{rk}\mathfrak u
  -\operatorname{rk}\operatorname{Centr}_{\mathfrak u}(X')
  \leq
  \dim U-\dim\operatorname{Centr}_U(X').
\]
Since \(U\) is smooth,
\[
  \dim U=\operatorname{rk}\mathfrak u.
\]
On the other hand, \SGARefLink{sga3:ii:paragraph:5-3-3}{Exposé II, 5.3.3} gives
\[
  \dim\operatorname{Centr}_U(X')
  \leq
  \operatorname{rk}\operatorname{Centr}_{\mathfrak u}(X'),
\]
which proves the assertion.  Notice in fact that
\[
  \dim\operatorname{Centr}_U(X')
  =
  \operatorname{rk}\operatorname{Centr}_{\mathfrak u}(X');
\]
this also proves once more that \(Z\) is smooth, as in \SGARefLink{sga3:xvii:lemma:5:7:1}{5.7.1}.
